\documentclass[12pt,a4paper]{article}

% ------------------------------------------------------------
% Pacotes
% ------------------------------------------------------------
\usepackage[utf8]{inputenc}
\usepackage[T1]{fontenc}
\usepackage[brazilian]{babel}
\usepackage[margin=2.5cm]{geometry}
\usepackage{graphicx}
\usepackage{enumitem}
\usepackage{listings}
\usepackage{booktabs}
\usepackage{tabularx}
\usepackage{xcolor}
\usepackage{float}    % permite usar [H] para fixar tabelas no lugar
\usepackage[hidelinks]{hyperref}

\setlist[itemize]{itemsep=2pt, topsep=4pt}

% Estilo dos trechos de código
\lstset{
  language=SQL,
  basicstyle=\ttfamily\footnotesize,
  frame=single,
  breaklines=true,
  breakatwhitespace=false,
  columns=fullflexible,
  keywordstyle=\color{blue},
  commentstyle=\color{gray},
  showstringspaces=false,
  xleftmargin=0pt,
  aboveskip=1em,
  belowskip=1em
}

% ------------------------------------------------------------
% PREENCHA AQUI: dados da capa
% ------------------------------------------------------------
\newcommand{\universidade}{Universidade Federal do Rio Grande -- FURG}
\newcommand{\disciplina}{Sistemas para Internet II}
\newcommand{\titulo}{Banco de dados}
\newcommand{\cidade}{Rio Grande}
\newcommand{\datacapa}{outubro de 2026}   % data escrita em português

% Textos em português (garantem português mesmo se o babel falhar)
\renewcommand{\figurename}{Figura}
\renewcommand{\tablename}{Tabela}
\renewcommand{\refname}{Referências}

\begin{document}

% ------------------------------------------------------------
% Capa
% ------------------------------------------------------------
\begin{titlepage}
  \centering
  {\Large \textbf{\universidade}\par}
  \vspace{1cm}
  {\large \disciplina\par}
  \vspace{3cm}
  {\LARGE \textbf{\titulo}\par}
  \vspace{3cm}
  {\large \textbf{Alunos:}\par}
  \vspace{0.4cm}
  {\large
    Lincon Avila de Souza -- Matrícula 162571\par
    Lucas Derick Silva Pereira -- Matrícula 162547\par
    Miguel Casarin da Silva -- Matrícula 162575\par}
  \vfill
  {\large \cidade, \datacapa\par}
\end{titlepage}

% O conteúdo do relatório começa aqui (sem sumário)

\section{Dados e Obtenção}

O banco de dados do projeto foi populado utilizando informações públicas extraídas do Portal de Dados Abertos, focando nos registros de servidores federais ativos (carreiras e cargos) e aposentados do Poder Executivo Federal Civil. Os dados brutos são disponibilizados em arquivos de texto e planilhas estruturadas, acompanhados por dicionários oficiais que descrevem detalhadamente os campos disponíveis.

\subsection{Dataset de Servidores Aposentados}
O arquivo voltado aos aposentados contempla o cadastro de inativos sob o regime jurídico único federal. Os principais campos extraídos e tratados a partir do dicionário original englobam:
\begin{itemize}
    \item \textbf{Nome}: Nome completo do servidor aposentado.
    \item \textbf{CPF}: Número de inscrição na Receita Federal, protegido por máscara de segurança.
    \item \textbf{Matrícula}: Identificador numérico exclusivo da matrícula do servidor.
    \item \textbf{Órgão e Sigla}: Nome e sigla do órgão público federal responsável pela dotação orçamentária vinculada.
    \item \textbf{Cargo}: Denominação do cargo efetivo que deu origem à aposentadoria.
    \item \textbf{Tipos e Legislação}: Informações complementares referentes ao tipo de aposentadoria, fundamentação legal e dados do diploma de publicação.
    \item \textbf{Ingresso e Proventos}: Data de ingresso no serviço público e valor líquido correspondente ao provento recebido.
\end{itemize}

\subsection{Dataset de Servidores Ativos (Carreiras e Cargos)}
O arquivo correspondente aos servidores ativos descreve o quadro de pessoal em atividade no âmbito do Poder Executivo Federal. Os campos essenciais identificados no dicionário de dados compreendem:
\begin{itemize}
    \item \textbf{Nome}: Nome completo do servidor ou empregado público.
    \item \textbf{CPF}: Número de cadastro de pessoa física mascarado.
    \item \textbf{Código da Carreira}: Código identificador do agrupamento profissional.
    \item \textbf{Descrição do Cargo}: Nome do cargo ou emprego público ocupado no plano de carreiras.
    \item \textbf{UF da UPAG}: Sigla da Unidade Federativa da Unidade Pagadora de vinculação.
    \item \textbf{Órgão de Atuação}: Nome do órgão orçamentário ao qual o trabalhador está vinculado.
    \item \textbf{Remuneração}: Valor líquido percebido pelo servidor durante o mês de referência.
\end{itemize}

\section{Exploração e testes na base de dados}

A fim de buscar possíveis candidatos a identificadores únicos, foram realizados testes com os arquivos \textit{CSV} sem tratamento. Identificou-se que nenhum campo era único na base de dados, havendo repetições, além de a \textbf{UF} não estar disponível no \textit{dataset} de aposentados. Outro ponto importante foi a observação de que o campo \textbf{CPF} possui seus dígitos identificadores substituídos por asteriscos, o que impediu o uso do \textbf{CPF} como chave entre as tabelas.

Os testes foram conduzidos através de \textit{queries} SQL, carregando os arquivos \textit{CSV} por meio do \textit{DuckDB} (banco de dados simples, que pode ser usado pelo terminal e cujo uso é focado em análise de dados). Foram usadas \textit{queries} como a seguinte para verificar repetição em colunas nas tabelas:

\begin{lstlisting}
SELECT nome, COUNT(*) AS qtd
FROM read_csv_auto('CARREIRA_072026.csv')
GROUP BY nome
HAVING COUNT(*) > 1
ORDER BY qtd DESC;
\end{lstlisting}

\section{Tratamento dos dados}
O tratamento dos dados foi realizado por meio dos \textit{scripts} escritos em \textit{Python} \textbf{\texttt{clear\_aposentados.py}} e \textbf{\texttt{clear\_carreira.py}}, disponíveis no diretório \textbf{\texttt{./db/etl}}.

Estes \textit{scripts} executam a limpeza e padronização dos arquivos brutos para garantir a compatibilidade com o SGBD. Em suma, as rotinas realizam as seguintes etapas:
\begin{itemize}
    \item \textbf{Remoção de caracteres inválidos}: O \textit{script} de aposentados limpa \textit{bytes} nulos e ajusta a codificação de caracteres (\textit{latin1} para \textit{UTF-8}) para evitar falhas de leitura.
    \item \textbf{Padronização de valores monetários}: Ambos os códigos convertem campos de remuneração e proventos do padrão brasileiro (com vírgulas para decimais e pontos para milhares) para o formato numérico padrão de banco de dados (ponto decimal).
    \item \textbf{Normalização de datas}: Conversão de campos textuais e de data para o formato padrão ISO (\texttt{YYYY-MM-DD}). No caso da carreira, a data de referência é gerada dinamicamente com base no mês informado e no ano extraído do nome do arquivo.
    \item \textbf{Garantia de integridade e estruturação}: Os dados limpos são salvos em novos arquivos delimitados por vírgulas (\textit{CSV}), padronizados com os cabeçalhos esperados pelas tabelas do banco de dados.
\end{itemize}

\section{Dicionário de Dados}

O projeto optou por não utilizar todas as colunas contidas nos arquivos do Portal de Dados Abertos. Ao todo, o \textit{dataset} de servidores ativos obteve um volume de 1.090.896 registros após a limpeza, enquanto o \textit{dataset} de aposentados obteve 677 registros, distribuídos nos seguintes campos:

\begin{table}[H]
\centering
\caption{Dicionário de dados da tabela \texttt{servidores\_ativos}.}
\label{tab:dicionario_ativos}
\small
\setlength{\tabcolsep}{4pt}
\begin{tabularx}{\textwidth}{l l l X}
\toprule
\textbf{Coluna (Schema Final)} & \textbf{Tipo SQL} & \textbf{Campo Original} & \textbf{Descrição / Formato} \\
\midrule
\texttt{id} & BIGINT PK & -- & Chave substituta (\textit{surrogate key}) \\
\texttt{nome} & VARCHAR(255) & Nome & Nome do servidor/empregado público \\
\texttt{cpf} & VARCHAR(11) & CPF & CPF mascarado (\texttt{***NNNNNN**}) \\
\texttt{codigo\_carreira} & VARCHAR(3) & Código da Carreira & Código de classes da profissão \\
\texttt{cargo} & VARCHAR(45) & Descrição do Cargo & Nome do cargo ou emprego ocupado \\
\texttt{uf} & VARCHAR(2) & UF da UPAG & Sigla da UF da unidade pagadora \\
\texttt{orgao} & VARCHAR(45) & Denominação do Órgão & Órgão federal de vínculo \\
\texttt{remuneracao} & NUMERIC(20,2) & Valor da Remuneração & Rendimento líquido recebido \\
\bottomrule
\end{tabularx}
\end{table}

\begin{table}[H]
\centering
\caption{Dicionário de dados da tabela \texttt{servidores\_aposentados}.}
\label{tab:dicionario_aposentados}
\small
\setlength{\tabcolsep}{4pt}
\begin{tabularx}{\textwidth}{l l l X}
\toprule
\textbf{Coluna (Schema Final)} & \textbf{Tipo SQL} & \textbf{Campo Original} & \textbf{Descrição / Formato} \\
\midrule
\texttt{id} & BIGINT PK & -- & Chave substituta (\textit{surrogate key}) \\
\texttt{nome} & VARCHAR(255) & Nome & Nome completo do aposentado \\
\texttt{cpf} & VARCHAR(11) & CPF & CPF parcialmente mascarado \\
\texttt{matricula} & BIGINT & Matrícula do Servidor & Matrícula do servidor (\texttt{OOOOOMMMMMMM}) \\
\texttt{orgao} & VARCHAR(45) & Nome do Órgão & Nome do órgão de vínculo \\
\texttt{sigla\_orgao} & VARCHAR(255) & Sigla do Órgão & Sigla institucional do órgão \\
\texttt{cargo} & VARCHAR(45) & Cargo & Cargo originário da aposentadoria \\
\texttt{tipo\_aposentadoria} & VARCHAR(255) & Tipo de Aposentadoria & Classificação sistêmica do tipo \\
\texttt{data\_aposentadoria} & DATE & Data publicação Diploma & Data de publicação legal \\
\texttt{tipo\_ingresso} & VARCHAR(255) & Ocorrência de ingresso & Tipo de ingresso no serviço público \\
\texttt{data\_ingresso} & DATE & Data ocorrência ingresso & Data efetiva de ingresso \\
\texttt{remuneracao} & NUMERIC(20,2) & Valor do Rendimento Líquido & Valor do provento líquido \\
\bottomrule
\end{tabularx}
\end{table}

\clearpage

\section{Tabelas e índices}

Para a persistência dos dados e a garantia de desempenho nas consultas, o SGBD PostgreSQL foi estruturado com duas tabelas principais e independentes (\texttt{servidores\_ativos} e \texttt{servidores\_aposentados}), além de índices otimizados para buscas exatas, compostas e por similaridade.

\begin{lstlisting}
CREATE EXTENSION IF NOT EXISTS pg_trgm;

CREATE TABLE servidores_ativos (
    id BIGINT PRIMARY KEY,
    nome VARCHAR(255) NOT NULL,
    cpf VARCHAR(11) NOT NULL, -- CPF censurado; nao restrito como UNIQUE devido a duplicatas de origem
    codigo_carreira VARCHAR(3) NOT NULL,
    cargo VARCHAR(45) NOT NULL,
    uf VARCHAR(2) NOT NULL,
    orgao VARCHAR(45) NOT NULL,
    remuneracao NUMERIC(20,2) NOT NULL
);

CREATE TABLE servidores_aposentados (
    id BIGINT PRIMARY KEY,
    nome VARCHAR(255) NOT NULL,
    cpf VARCHAR(11) NOT NULL,
    matricula BIGINT NOT NULL,
    orgao VARCHAR(45) NOT NULL,
    sigla_orgao VARCHAR(255) NOT NULL,
    cargo VARCHAR(45) NOT NULL,
    tipo_aposentadoria VARCHAR(255) NOT NULL,
    data_aposentadoria date NOT NULL,
    tipo_ingresso VARCHAR(255) NOT NULL,
    data_ingresso date NOT NULL,
    remuneracao NUMERIC(20,2) NOT NULL
);

-- Indices B-tree para buscas exatas isoladas
CREATE INDEX IF NOT EXISTS idx_servidores_ativos_nome ON servidores_ativos USING BTREE (nome);
CREATE INDEX IF NOT EXISTS idx_servidores_ativos_cargo ON servidores_ativos USING BTREE (cargo);
CREATE INDEX IF NOT EXISTS idx_servidores_ativos_uf ON servidores_ativos USING BTREE (uf);
CREATE INDEX IF NOT EXISTS idx_servidores_ativos_orgao ON servidores_ativos USING BTREE (orgao);

CREATE INDEX IF NOT EXISTS idx_servidores_aposentados_nome ON servidores_aposentados USING BTREE (nome);
CREATE INDEX IF NOT EXISTS idx_servidores_aposentados_cargo ON servidores_aposentados USING BTREE (cargo);
CREATE INDEX IF NOT EXISTS idx_servidores_aposentados_orgao ON servidores_aposentados USING BTREE (orgao);

-- Indices em minusculo para otimizacao de busca
CREATE INDEX IF NOT EXISTS idx_servidores_ativos_nome_lower ON servidores_ativos USING BTREE (LOWER(nome));
CREATE INDEX IF NOT EXISTS idx_servidores_aposentados_nome_lower ON servidores_aposentados USING BTREE (LOWER(nome));

-- Indices compostos para multiplos filtros simultaneos
CREATE INDEX IF NOT EXISTS idx_servidores_ativos_cargo_uf ON servidores_ativos USING BTREE (cargo, uf);
CREATE INDEX IF NOT EXISTS idx_servidores_ativos_orgao_uf ON servidores_ativos USING BTREE (orgao, uf);

-- Indices GIN com pg_trgm para busca por similaridade 
CREATE INDEX IF NOT EXISTS idx_servidores_ativos_nome_trgm ON servidores_ativos USING GIN (nome gin_trgm_ops);
CREATE INDEX IF NOT EXISTS idx_servidores_aposentados_nome_trgm ON servidores_aposentados USING GIN (nome gin_trgm_ops);
\end{lstlisting}

\end{document}